\subsection{Sprint 0}

Com o início do projeto após a reunião com a stakeholder, identificamos problemas graves com
o Termo de Abertura: todos acreditávamos que seria um projeto mobile simples, voltado a uma
\ac{api} específica para esse ambiente. No entanto, durante a conversa com a stakeholder, ela
explicou que o que estava descrito no termo era apenas um pequeno pedaço de um \textit{hub}
--- um ecossistema da ENSportive com diversas aplicações ligadas a ele, como mentorias,
cursos de saúde, formação de professores, empresas e colaboradores, e a escola de raquetes,
entre outros.

Com um posicionamento muito assertivo do nosso colega de \ac{ages} IV, conseguimos delimitar o
escopo exclusivamente para a Escola de Raquetes, visto que não haveria tempo hábil para
construir o ecossistema completo. Definiu-se que, futuramente, nossa aplicação poderia ser
integrada como um dos microsserviços desse hub. Com a proposta definida e aceita pela
stakeholder, iniciou-se a fase de definições técnicas.

Identificamos as ferramentas que seriam utilizadas: \textbf{Spring Boot} para o backend,
com \textbf{Maven} como gerenciador de dependências. Para o frontend, o time se dividiu:
parte optou por ferramentas já conhecidas, enquanto outra parte decidiu aprender
\textbf{Angular 17}, uma versão que trazia novidades significativas e nos faria sair do
padrão React, mais comum entre os alunos de \ac{ages}. Após muitas discussões, o grupo foi
convencido por um colega que trabalhava com Angular 17 profissionalmente, e a decisão foi
unânime.

Os alunos de \ac{ages} II ficaram responsáveis pelos \textit{mockups}, enquanto os de \ac{ages} I
se dedicaram ao estudo das ferramentas. Os \ac{ages} III cuidaram da estruturação do Git e dos
requisitos, junto com os \ac{ages} IV. De forma colaborativa, todos contribuíram na busca de
bibliotecas funcionais para o calendário e na referência de exemplos de interface.
Os mockups foram concluídos dentro do prazo previsto para apresentação à stakeholder.